% Part: computability
% Chapter: computability-theory
% Section: def-functions-self-reference

\documentclass[../../../include/open-logic-section]{subfiles}

\begin{document}

\olfileid{cmp}{thy}{slf}
\olsection{تعریف توابع با خودارجاعی}

به‌طور کلی، توانایی تعریف توابع برحسب خودشان سودمند است. برای نمونه،
با داشتن توابع محاسبه‌پذیر $k$، $l$ و~$m$، لم نقطهٔ ثابت می‌گوید
تابع محاسبه‌پذیر جزئیِ~$f$ای وجود دارد که به ازای هر~$y$ معادلهٔ زیر
را برآورده می‌کند:
\[
f(y) \simeq
\begin{cases}
k(y) & \text{اگر $l(y) = 0$} \\
f(m(y)) & \text{در غیر این صورت.}
\end{cases}
\]
به‌طور مشخص‌تر، $f$ با قراردادن
\[
g(x,y) \simeq
\begin{cases}
k(y) & \text{اگر $l(y) = 0$} \\
\cfind{x}(m(y)) & \text{در غیر این صورت}
\end{cases}
\]
و سپس به‌کاربردن لم نقطهٔ ثابت برای یافتن نمایه‌ای چون~$e$ به دست
می‌آید که $\cfind{e}(y)\simeq g(e,y)$.

به‌عنوان نمونه‌ای عینی، تابع «بزرگ‌ترین مقسوم‌علیه مشترک»
$\fn{gcd}(u,v)$ را می‌توان چنین تعریف کرد:
\[
\fn{gcd}(u,v) \simeq
\begin{cases}
v & \text{اگر $u = 0$} \\
\fn{gcd}(\fn{mod}(v, u), u) & \text{در غیر این صورت}
\end{cases}
\]
که در آن $\fn{mod}(v,u)$ باقیماندهٔ تقسیم $v$ بر~$u$ را نشان می‌دهد.
استناد به لم نقطهٔ ثابت نشان می‌دهد $\fn{gcd}$ محاسبه‌پذیر جزئی است.
(در واقع می‌توان این را با قراردادن کد زوج $\tuple{u,v}$ در~$y$، در
قالب بالا نوشت.) سپس استقرا بر~$u$ نشان می‌دهد $\fn{gcd}$ در واقع
تام است.

البته می‌توان تعریف‌های خودارجاعیِ بسیار پیچیده‌تری از نمونه‌های
بررسی‌شده ساخت. بیشتر زبان‌های برنامه‌نویسی به شکلی از تعریف توابع
برحسب خودشان پشتیبانی می‌کنند. توجه کنید این قدری کم‌هیجان‌تر از
توانایی تعریف تابع برحسب \emph{نمایهٔ} الگوریتمی است که تابع را محاسبه
می‌کند؛ امکانی که قضیهٔ نقطهٔ ثابت با کلیت کامل به دست می‌دهد.

\end{document}
